\documentclass[12pt]{article}

% --- Packages ---
\usepackage[utf8]{inputenc}
\usepackage[T1]{fontenc}
\usepackage[french]{babel}
\usepackage{amsmath, amssymb}
\usepackage{geometry}
\usepackage{tikz}
\usepackage[most]{tcolorbox}
\usepackage{multicol}
\usepackage{lmodern}
\usepackage{pgfplots}
\pgfplotsset{compat=1.18}
\usepackage{mdframed}
\usepackage{bm}
\usepackage{xcolor}
\usepackage{sectsty}
\usepackage[utf8]{inputenc}
\usepackage[T1]{fontenc}
\usepackage[inline]{enumitem}
\usepackage[dvipsnames]{xcolor}
\usepackage{tkz-tab}
% --- Page ---
\geometry{margin=2cm}

% Footer
\usepackage{fancyhdr}
\pagestyle{fancy}
\cfoot{\thepage}
\rfoot{\textit{G.M.}}
\renewcommand{\headrulewidth}{0pt}
\renewcommand{\footrulewidth}{0pt}

% Pour les normes : 
\newcommand*{\norme}[1]{\left\lVert\overrightarrow{#1}\right\rVert}

\usetikzlibrary{shapes.misc}

\tikzset{cross/.style={cross out, draw=black, minimum size=2*(#1-\pgflinewidth), inner sep=0pt, outer sep=0pt},
%default radius will be 1pt. 
cross/.default={4pt}}

% mettre dans le préambule, après \usepackage{xcolor}
\usepackage{sectsty} % gère facilement la couleur des titres
\setlength{\parindent}{0pt}

\sectionfont{\color{red}\Large\bfseries}     % pour \section*
% \subsectionfont{\color{red}\normalsize\bfseries} % pour \subsection*

\begin{document}

% ============================
%   ÉQUATIONS
% ============================

\begin{center}
    {\LARGE \textbf{Équations}}\\[4mm]
    {\large Seconde}
\end{center}

\hrule
\vspace{0.4cm}

% ============================
%   PARTIE 1 : Équation du premier degré
% ============================

\section*{\textcolor{red}{Partie 1 : Équation du premier degré}}

\subsection*{\textcolor{red}{1.1 Définition}}

\vspace{5pt}

\begin{mdframed}[
    linewidth=1pt,
    % roundcorner=5pt,
    linecolor=Red,
    innertopmargin=10pt,
    % innerbottommargin=10pt
]
{\color{red}
\textbf{\underline{But :} Trouver x !}
\vspace{5pt}\\
C'est-à-dire isoler $x$ dans l'équation pour arriver à :
\begin{center}
    $x=$ nombre
\end{center}
\vspace{0pt}
}
\end{mdframed}

\vspace{10pt}

\underline{Exemple :}
\[
3x - 7 = 5
\]

\vspace{5pt}

\underline{Résolution :}
\[
3x - 7 \textcolor{red}{\ +\ 7} = 5 \textcolor{red}{\ +\ 7}
\]
\[
3x = 12
\]
\[
\dfrac{3x}{\textcolor{red}{3}} = \dfrac{12}{\textcolor{red}{3}}
\]
\[
x = 4
\]

\vspace{10pt}

\begin{mdframed}[
    linewidth=1pt,
    roundcorner=5pt,
    linecolor=Red
]
{\color{red}
\underline{Méthode :}
\begin{itemize}
    \item On regroupe les termes contenant $x$ d'un côté.
    \item On regroupe les nombres de l'autre côté.
\end{itemize}
}
\end{mdframed}

\vspace{10pt}

\begin{mdframed}[
    linewidth=1pt,
    roundcorner=5pt,
    linecolor=RoyalBlue,
    innertopmargin=10pt,
    innerbottommargin=10pt
]
{\color{RoyalBlue}\textbf{\underline{Exercice d'application n°1 :}}}
\vspace{5pt}
\\
Résoudre les équations suivantes :
\begin{enumerate}
    \item $7x - 4 = 3x + 8$
    \item $\dfrac{4}{3}x - 2 = 7 - \dfrac{5}{2}x$
    \item $3(x+4) = -(x+5) + 2$
\end{enumerate}
\end{mdframed}

\newpage

% ============================
%   PARTIE 2 : Équation produit
% ============================

\section*{\textcolor{red}{Partie 2 : Équation produit}}

\subsection*{\textcolor{red}{2.1 Définition}}

\vspace{5pt}

\begin{mdframed}[
    linewidth=1pt,
    roundcorner=5pt,
    linecolor=ForestGreen
]
{\color{ForestGreen}\textbf{Définition :}}  
Une \textbf{équation produit} est une équation de la forme :
\[
P(x) \times Q(x) = 0
\]
Où $P(x)$ et $Q(x)$ sont des expressions littérales.
\end{mdframed}

\vspace{10pt}

\begin{mdframed}[
    linewidth=1pt,
    roundcorner=5pt,
    linecolor=Red,
    innertopmargin=10pt,
    % innerbottommargin=10pt
]
{\color{Red}\textbf{\underline{Propriété fondamentale :}}
\begin{center}
    Si $P(x) \times Q(x) = 0$ alors  $P(x)=0 \ \text{ou} \ Q(x)=0$
\end{center}
}
\end{mdframed}

\underline{Autre formulation :}
\vspace{5pt}
\\Si un produit de facteurs est nul, alors 
au moins l'un des facteurs \_\_\_\_\_\_\_\_\_\_\_\_.

\vspace{10pt}

\underline{Exemple :}
\[
(x - 3)(2x + 1) = 0
\]

\vspace{5pt}

\underline{Résolution :}
\[
x - 3 = 0 \quad \text{\underline{ou}} \quad 2x + 1 = 0
\]
\[
x = 3 \quad \text{\underline{ou}} \quad x = -\frac{1}{2}
\]

\vspace{10pt}

\begin{mdframed}[
    linewidth=1pt,
    roundcorner=5pt,
    linecolor=RoyalBlue,
    innertopmargin=10pt,
    innerbottommargin=10pt
]
{\color{RoyalBlue}\textbf{\underline{Exercice d'application n°2 :}}}
\vspace{5pt}
\\
Résoudre les équations suivantes :
\begin{enumerate}
    \item $(4x + 6)(3 - 7x) = 0$
    \item $(3x + 1)(1 - 6x) - (3x + 7)(3x + 1) = 0$
    \item $5x^2 - 4x = 0$
\end{enumerate}
\end{mdframed}

% ============================
%   PARTIE 3 : Équation de la forme x² = a
% ============================

\section*{\textcolor{red}{Partie 3 : Équation de la forme $x^2 = a$}}

\vspace{5pt}

\begin{mdframed}[
    linewidth=1pt,
    roundcorner=5pt,
    linecolor=ForestGreen
]
{\color{ForestGreen}\textbf{Définition :}}  
On considère une équation de la forme :
\[
x^2 = a
\]
où $a$ est un nombre réel.
\end{mdframed}

\vspace{10pt}

\begin{mdframed}[
    linewidth=1pt,
    roundcorner=5pt,
    linecolor=Red,
    innertopmargin=10pt,
    innerbottommargin=10pt
]
{\color{red}\underline{\textbf{Propriété :}}
\vspace{5pt}\\
Les solutions dans $\mathbb{R}$ de l'équation $x^2 = a$ dépendent du signe de $a$ :
\begin{itemize}
    \item Si $a > 0$, alors l'équation admet \textbf{deux solutions} :  
    $x = \sqrt{a}$ et $x = -\sqrt{a}$
    \item Si $a = 0$, alors l'équation possède une \textbf{unique solution} : $x = 0$
    \item Si $a < 0$, alors l'équation n'a \textbf{aucune solution}
\end{itemize}
}
\end{mdframed}

\vspace{10pt}

\underline{Exemple :}
\[
x^2 = 9
\]

\vspace{5pt}

\underline{Résolution :}
\[
x = 3 \quad \text{ou} \quad x = -3
\]

\vspace{10pt}

\begin{mdframed}[
    linewidth=1pt,
    roundcorner=5pt,
    linecolor=RoyalBlue,
    innertopmargin=10pt,
    innerbottommargin=10pt
]
{\color{RoyalBlue}\textbf{\underline{Exercice d'application n°3 :}}}
\vspace{5pt}
\\
Résoudre les équations suivantes :
\begin{enumerate}
    \item $x^2 = 16$
    \item $x^2 = -8$
    \item $(x+2)^2 = 9$
\end{enumerate}
\end{mdframed}

% ============================
%   PARTIE 4 : Équation quotient
% ============================

\section*{\textcolor{red}{Partie 4 : Équation quotient}}

\vspace{5pt}

\begin{mdframed}[
    linewidth=1pt,
    roundcorner=5pt,
    linecolor=ForestGreen
]
{\color{ForestGreen}\textbf{Définition :}}  
Une \textbf{équation quotient} est une équation de la forme :
\[
\frac{P(x)}{Q(x)} = 0
\]
Où $P(x)$ et $Q(x)$ sont des expressions littérales ($\bm{Q(x) \neq 0}$).
\end{mdframed}

\vspace{10pt}

\begin{mdframed}[
    linewidth=1pt,
    roundcorner=5pt,
    linecolor=Red,
    innertopmargin=10pt,
    % innerbottommargin=10pt
]
{\color{red}\underline{Propriété :}
\begin{center}
    Si $\dfrac{P(x)}{Q(x)} = 0$ alors $P(x) = 0$ \textbf{\underline{ET}} $Q(x) \neq 0$.
\end{center}
}
\end{mdframed}

\vspace{10pt}

\underline{Exemple :}
\[
\frac{2x - 4}{x + 1} = 0
\]

\vspace{5pt}

\underline{Résolution :}
\[
2x - 4 = 0 \quad \text{et} \quad x + 1 \neq 0
\]
\[
2x = 4 \quad \text{et} \quad x \neq -1
\]
\[
x = 2 \quad \text{et} \quad x \neq -1
\]
\begin{center}
    $x = 2$ est une solution valide car elle n'est pas égale à $-1$ (valeur interdite)
\end{center}

\vspace{10pt}

\begin{mdframed}[
    linewidth=1pt,
    roundcorner=5pt,
    linecolor=RoyalBlue,
    innertopmargin=10pt,
    innerbottommargin=10pt
]
{\color{RoyalBlue}\textbf{\underline{Exercice d'application n°4 :}}}
\vspace{5pt}
\\
Résoudre dans $\mathbb{R}$ les équations suivantes :
\begin{multicols}{2}
    \begin{enumerate}
        \item $\dfrac{3x+5}{x-1} = 0$
        \vspace{5pt}
        \item $\dfrac{(2x+1)(x-3)}{x-4} = 0$
        \vspace{5pt}
        \item $\dfrac{x^2 - 9}{x+3} = 0$
        \vspace{5pt}
        \item $\dfrac{x+3}{x-3} = \dfrac{2}{x-3}$
        \vspace{5pt}
        \item (difficile) $1 - \dfrac{x+3}{x-3} = \dfrac{2}{2-x}$
    \end{enumerate}
\end{multicols}
\end{mdframed}

\end{document}
